% 品性与修养.tex
% Auto-generated from 感悟论语.odt by tools/odt_to_tex.py

\setpartplant{artwork/pine.png}
  {松柏}
  {岁寒，然后知松柏之后凋也。}
  {《论语·子罕》}
  {品性与修养}
\part{品性与修养}

\chapter{君子品性与儒家精神}

君子不拘泥于表面的嘘寒问暖，他们关注的是在追求至善道路上的慎思笃行。君子自强不息、思想与精神独立，靠的不是相濡以沫、弹冠相庆，而是相望于江湖、各守善道，享一份“有朋自远方来”的乐处便好。尽管“人不己知”，亦不会有愠色，甚至还能激发出“患其不能”的无穷动力。因为深切明白“可与共学，未可与适道；可与适道，未可与立；可与立，未可与权”的道理，即便知音难觅，鲜有人可与碰撞出思想的火花，也同样不忧不惧。君子不担心饮食是否精美、衣着是否时尚、出行是否排场，因为一箪食、一瓢饮便足以支持其终生学习和对真理的追求。由这样的品性，便自然联想到我先前总结的儒家精神。在此重提，以警醒自我：

真正的儒家，有止于至善的终极目标，有君子固穷的气节。既有危邦不入、乱邦不居的小心慬慎，又有可以仕而仕的智慧敏锐。能于箪食瓢饮之间体味人生乐趣，亦有风乎舞雩、歌咏而归的豁达情怀。因为心怀至善而宽容，因为有原则而不妥协。终生学习，以博学、明辨、审问、笃行为乐，以巧言令色、浅尝辄止为耻。

\begin{classic}
《论语·学而篇第一 1.1》 子曰：「学而时习之，不亦说乎？有朋自远方来，不亦乐乎？人不知而不愠，不亦君子乎？」

《论语·宪问篇第十四 14.30》 子曰：「不患人之不己知，患其不能也。」

《论语·子罕篇第九 9.30》 子曰：「可与共学，未可与适道；可与适道，未可与立；可与立，未可与权。」

《论语·颜渊篇第十二 12.4》 司马牛问君子。子曰：「君子不忧不惧。」曰：「不忧不惧，斯谓之君子矣乎？」子曰：「内省不疚，夫何忧何惧？」

《论语·雍也篇第六 6.11》 子曰：「贤哉，回也！一箪食，一瓢饮，在陋巷，人不堪其忧，回也不改其乐。贤哉，回也！」

《论语·卫灵公篇第十五 15.2》 在陈绝粮。从者病，莫能兴。子路愠见曰：「君子亦有穷乎？」子曰：「君子固穷，小人穷斯滥矣。」

《论语·泰伯篇第八 8.13》 子曰：「笃信好学，守死善道。危邦不入，乱邦不居，天下有道则见，无道则隐。邦有道，贫且贱焉，耻也，邦无道，富且贵焉，耻也。」

《论语·先进篇第十一 11.26》 子路、曾皙、冉有、公西华侍坐。子曰：「以吾一日长乎尔，毋吾以也。居则曰：「不吾知也！」如或知尔，则何以哉？」子路率尔而对曰：「千乘之国，摄乎大国之间，加之以师旅，因之以饥馑，由也为之，比及三年，可使有勇，且知方也。」夫子哂之。「求，尔何如？」对曰：「方六七十，如五六十，求也为之，比及三年，可使足民；如其礼乐，以俟君子。」「赤，尔何如？」对曰：「非曰能之，愿学焉！宗庙之事，如会同，端章甫，愿为小相焉。」「点，尔何如？」鼓瑟希，铿尔，舍瑟而作。对曰：「异乎三子者之撰。」子曰：「何伤乎？亦各言其志也。」曰：「莫春者，春服既成；冠者五六人，童子六七人，浴乎沂，风乎舞雩，咏而归。」夫子喟然叹曰：「吾与点也！」三子者出，曾皙后。曾皙曰：「夫三子者之言何如？」子曰：「亦各言其志也已矣！」曰：「夫子何哂由也？」曰：「为国以礼，其言不让，是故哂之。」「唯求则非邦也与？」「安见方六七十，如五六十，而非邦也者。」「唯赤，非邦也与？」「宗庙会同，非诸侯而何？赤也为之小，孰能为之大！」

《论语·颜渊篇第十二 12.15》 子曰：「博学于文，约之以礼，亦可以弗畔矣夫。」

《论语·阳货篇第十七 17.17》 子曰：「巧言令色鲜矣仁。」
\end{classic}

\quotsource{《礼记 儒行》}

\begin{classic}
儒有席上之珍以待聘 ，夙夜强学以待问 ，怀忠信以待举 ，力行以待取 。其自立有如此者 。

儒有可亲而不可劫也 ，可近而不可迫也 ，可杀而不可辱也 。其居处不淫 ，其饮食不溽 ，其过失可微辨而不可面数也 。其刚毅有如此者 。

儒有今人与居，古人与稽，今世行之，后世以为楷。适弗逢世，上弗援，下弗推，谗谄之民有比党而危之者，身可危也，而志不可夺也。虽危，起居竟信其志，犹将不忘百姓之病也。其忧思有如此者。

儒有博学而不穷，笃行而不倦，幽居而不淫，上通而不困。礼之以和为贵，忠信之美，优游之法，慕贤而容众，毁方而瓦合。其宽裕有如此者。

儒有闻善以相告也，见善以相示也，爵位相先也，患难相死也，久相待也，远相致也。其任举有如此者。

儒有澡身而浴德，陈言而伏，静而正之。上弗知也，麤而翘之，又不急为也。不临深而为高，不加少而为多。世治不轻，世乱不沮，同弗与，异弗非也。其特立独行有如此者。

温良者，仁之本也;敬慎者，仁之地也;宽裕者，仁之作也;孙接者，仁之能也;礼节者，仁之貌也;言谈者，仁之文也;歌乐者，仁之和也;分散者，仁之施也。儒皆兼此而有之，犹且不敢言仁也。其尊让有如此者。
\end{classic}

\entrydate{2017.07.27}

\chapter{梁启超《论君子》注释}

君子二字其意甚广，欲为之诠注，颇难得其确解。为英人所称劲德尔门（gentleman）包罗众义，与我国君子之意差相吻合。证之古史，君子每与小人对待【《论语·里仁第四 4.16》子曰：“君子喻于义，小人喻于利。”《论语·述而第七 7.37》子曰：“君子坦荡荡，小人长戚戚。”《论语·卫灵公篇第十五 15.2》 在陈绝粮。从者病，莫能兴。子路愠见曰：“君子亦有穷乎？”子曰：“君子固穷，小人穷斯滥矣。”】，学善则为君子【《论语·学而篇第一 1.1》 子曰：“学而时习之，不亦说乎？有朋自远方来，不亦乐乎？人不知而不愠，不亦君子乎？”】，学不善则为小人。君子小人之分，似无定衡。顾习尚沿传类以君子为人格之标准。望治者，每以人人有士君子之心相勖【勖xù：勉励】。《论语》云：“君子人与？君子人也！” 明乎君子品高，未易几及也。【《论语·泰伯篇第八 8.6》 曾子曰：「可以讬六尺之孤，可以寄百里之命，临大节而不可夺也，君子人与，君子人也。」】

英美教育精神，以养成国民之人格为宗旨。国家犹机器也，国民犹轮轴也。转移盘旋，端在国民，必使人人得发展其本能，人人得勉为劲德尔门，即我国所谓君子者。莽莽神州，需用君子人，於今益极，本英美教育大意而更张之。国民之人格，駸駸【駸qīn：马极速奔驰貌】日上乎。

君子之义，既鲜确诂，欲得其具体的条件，亦非易言。《鲁论》所述，多圣贤学养之渐【渐：条件，前提；逐步发展的过程】，君子立品之方，连篇累牍势难胪【胪lú：陈述】举。周易六十四卦，言君子者凡五十三。乾坤二卦所云尤为提要钧元。乾象曰：“天行健，君子以自强不息。”坤象曰：“地势坤，君子以厚德载物。”推本乎此，君子之条件庶几近之矣。

乾象言，君子自励犹天之运行不息，不得有一暴十寒之弊。才智如董子，犹云勉强学问。《中庸》亦曰“或勉强而行之”。人非上圣，其求学之道，非勉强不得入于自然。且学者立志，尤须坚忍强毅，虽遇颠沛流离【《论语·里仁篇第四 4.5》 子曰：“富与贵，是人之所欲也，不以其道得之，不处也。贫与贱，是人之恶也，不以其道得之，不去也。君子去仁，恶乎成名。君子无终食之间违仁，造次必于是，颠沛必于是。”】，不屈不挠；若或见利而进，知难而退，非大有为者之事，何足取焉？人之生世，犹舟之航于海，顺风逆风，因时而异，如必风顺而后扬帆，登岸无日矣。

且夫自胜则为强，乍见孺子入水，急欲援手，情之真也【《孟子·告子章句上》所以谓人皆有不忍人之心者，今人乍见孺子将入于井，皆有怵惕恻隐之心；】。继而思之，往援则己危，趋而避之，私欲之念起，不克自胜故也。孔子曰：“克己复礼为仁。”【《论语·颜渊篇第十二 12.1》 颜渊问仁。子曰：“克己复礼，为仁。一日克己复礼，天下归仁焉。为仁由己，而由仁乎哉？”颜渊曰：“请问其目？”子曰：“非礼勿视，非礼勿听，非礼勿言，非礼勿动。”颜渊曰：“回虽不敏，请事斯语矣！”】王阳明曰：“治山中贼易，治心中贼难。”古来忠臣孝子愤时忧国奋不欲生，然或念及妻儿，辄有难于一死不能自克者。若能摈私欲尚果毅，自强不息，则自励之功与天同德，犹英之劲德尔门，见义勇为【《论语·为政篇第二 2.24》子曰：「非其鬼而祭之，谄也。见义不为，无勇也。」】，不避艰险，非吾辈所谓君子其人哉？

坤象言，君子接物，度量宽厚，犹大地之博，无所不载。君子责己甚厚，责人甚轻。孔子曰：“躬自厚而薄责于人。”【《论语·卫灵公篇第十五 15.15》 子曰：「躬自厚，而薄责于人，则远怨矣！」】盖惟有容人之量，处世接物坦焉无所芥蒂，然后得以膺重任；非如小有才者，轻佻狂薄，毫无度量，不然小不忍必乱大谋，君子不为也。当其名高任重，气度雍容，望之俨然，即之温然【《论语·子张篇第十九 19.9》 子夏曰：「君子有三变：望之俨然，即之也温，听其言也厉。」】，此其所以为厚也，此其所以为君子也。

纵观四万万同胞，得安居乐业，教养其子若弟者几何人？读书子弟能得良师益友之薰陶者几何人？清华学子，荟中西之鸿儒，集四方之俊秀，为师为友，相蹉相磨，他年遨游海外，吸收新文明，改良我社会，促进我政治。所谓君子人者，非清华学子，行将焉属？虽然君子之德风，小人之德草【《论语·颜渊篇第十二 12.19》 季康子问政于孔子曰：“如杀无道，以就有道，何如？”孔子对曰：“子为政，焉用杀？子欲善，而民善矣！君子之德风，小人之德草，草上之风必偃。”】，今日之清华学子，将来即为社会之表率，语、默、作、止，皆为国民所仿效。设或不慎，坏习惯之传行急如暴雨，则大事偾【偾fèn：覆败，灭亡】矣。深愿及此时机，崇德【《论语·颜渊篇第十二 12.21》 樊迟从游于舞雩之下。曰：“敢问崇德、修慝、辨惑？”子曰：“善哉问！先事后得，非崇德与？攻其恶，无攻人之恶，非修慝与？一朝之忿，忘其身以及其亲，非惑与？”】修学【《论语·述而篇第七 7.3》 子曰：“德之不修，学之不讲，闻义不能徒，不善不能改，是吾忧也。”】，勉为真君子，异日出膺大任，足以挽既倒之狂澜，作中流之底柱，则民国幸甚矣。

\entrydate{2020.10.11}

\chapter{己之所欲，勿施于人}

一个人走入社会的时间愈久，就会发现人们为了一己之利，使是非黑白之界线变得愈不分明。一直以来自认为正确与高尚的原则，也会被许多人以不识时务为由而嗤之以鼻。求得别人的理解与认可已成为不可能，能够不被误解就已实属不易了。这时，自己若没有经过理性思考后属于自己的一份追求，便会要么困惑，要么麻木。然后，认可和误解都沦为表象，剩下的不过是见面时的寒暄，酒桌上的投缘，分别后的冷漠。但若能够始终秉持自己的纯真原则与至善之地，则认可与误解就都会变得轻如鸿毛，因为真理与美好本身并不会因为认可或误解而有所盈亏。正所谓“不曰坚乎，磨而不磷。不曰白乎，涅而不缁”。既然如此，期望与人交流、相互理解以至成为知己，就要看机缘，不宜强求。对于持有不同观点乃至不同人生观、世界观、价值观的人，其实根本没有什么可谈和可交流的，也就更不必去尝试将其说服和改变了。否则，对牛弹琴，自取辱焉。

春秋战国时期，孔子的孙子子思，《中庸》著者，是鲁国国君缪公的老师。缪公想要与子思成为朋友。子思说，论地位，你是国君，我是臣子，我们不可能成为朋友。而在学识与德行修养上，你是我的学生，我们又怎么可能成为朋友？（见于《孟子·万章下》：缪公亟见于子思，曰：“古千乘之国以友士，何如？”子思不悦，曰：“古之人有言曰，事之云乎，岂曰友之云乎？”子思之不悦也，岂不曰：“以位，则子，君也；我，臣也；何敢与君友也？以德，则子事我者也，奚可以与我友？”）所以，子曰“无友不如己者”，或者更准确地说，只与志同道合的人交友，是也。

已如上面所说，无需和不同道的人说太多，但那些有一定学识与修养的人，却往往情不自禁地向身边的人宣讲自己认为正确的真理和原则。虽说这也是怀着一颗好心，但若选错了对象，就很可能反而招来是非。比如，你对于自己的作品一向有高质量甚至是近乎苛刻的要求。你自觉这是再自然与正确不过的原则——事实也的确如此。但你若将此标准施于他人，甚至是身边整个浮躁、粗心、吹嘘、造假的环境，这无异于断了他人赖以生存之道——即便这样的生存之道在你自己看来是那么地低俗和肮脏，但这毕竟是他人的生存之道啊！孔子云“己所不欲，勿施于人”。现在看来，还需要再加上一句，己之所欲，亦勿施于人。这样一来，便可以将自己从“对牛弹琴”的尴尬局面中解脱出来，把关注点真正地聚焦在自己身上——首先明确属于独特的自己的那个至善至地，然后笃实践行，而不是在自己还做不到、做得不够、理想之实现仍旧遥遥无期的时候便迫不及待地想要说服人、感化人。这时的自己还不具备这样的资格和影响力。当然，此时若能遇上知己一二，可以相互理解彼此追求的意义、求索的艰难、其中精神的可贵，则是人生一幸事。但若无此知己，则亦当独自勇往前行，义无返顾。

以上推理与阐述似乎有把本来当是快乐轻松的人生搞得无比悲惨、孤苦之嫌，且大有从此天涯孤旅之意。但孔子亦有云，“求仁得仁又何怨”。所以，悲苦之极当是快乐与幸福之极。人生也正是无处不充满如此二律悖反的哲学涵义。更何况，若从整个家族传承或基因的沿续来看，对于处于一般家庭的人而言，能依照《大学》篇明确自己的“至善之地”，然后终生奋斗“止于至善”，恐怕其本人便是始作俑者。这就像是富贵之家于事业草创之时无比艰难，在理性与唯美世界中的开拓者，亦不会轻松。然而，“积善之家，必有馀庆”，在正确的方向上笃行、积累，获益的必不止其自身，更是为其后代打下好的基础。即便是在自己的有生之年无法影响他人，其精神或可由后代予以传承和发扬。站在更长的时间线上去思考此问题，岂不是更能有所宽慰？

\begin{classic}
《论语·颜渊篇第十二 12.2》 仲弓问仁。子曰：「出门如见大宾，使民如承大祭。己所不欲，勿施于人。在邦无怨，在家无怨。」仲弓曰：「雍虽不敏，请事斯语。」

《论语·阳货篇第十七 17.7》 佛肸召，子欲往。子路曰：「昔者由也闻诸夫子曰：『亲于其身为不善者，君子不入也』。佛肸以中牟畔，子之往也，如之何？」子曰：「然，有是言也。不曰坚乎？磨而不磷；不曰白乎？涅而不缁。吾岂匏瓜也哉？焉能系而不食！」

《论语·学而篇第一 1.8》 子曰：「君子不重，则不威，学则不固。主忠信，无友不如己者，过则勿惮改。」
\end{classic}

\entrydate{2016.08.28}

\chapter{刚毅木讷近仁}

没有深厚的积淀，没有经过千百次的提纯与锤炼，即便是费尽心思、绞尽脑汁写出来的东西也没有什么实际的价值，更不用提深远的影响力了。不仅如此，无论是表面上还是潜在的浮躁就像是一副镣铐，牢牢地将个人的水平钳制和锁死在低级的层次上，断绝了能够进一步提升的机会。须知，这和学术文章的发表是一个道理：凑数终究容易，上层次实为困难。这个“难”不只是“上层次”这个目标本身的困难，还要直面在（看似）毫无建树的漫长积累过程中来自他人的质疑、鄙夷、讥笑，以及当自我信念不够坚定、内心不够强大时，来自于内心和灵魂深处的怀疑犹豫、摇摆不定。

随着年龄的不断增长，一个人应当从外向型到内向型逐步地收缩、回退，从而变得沉稳、坚实；亦应抛开世俗成败与名利的困扰，回归到人性向善的本心；还要依据事物正常而健康发展所应遵循的基本规律谨慎、笃实行事，并以此为乐、为荣。这样一来，人就不会轻易地像公知们那样发表言论和文章了。更何况，说出去的漂亮话、写出来的严谨文字，自己能真正做到多少，又能收获几分的心得体会呢？刚毅木讷近仁是也！

\begin{classic}
《论语·子路篇第十三 13.27》 子曰：「刚毅木讷，近仁。」
\end{classic}

\entrydate{2018.11.05}

\chapter{由执鞭之士到最好的自己}

人一旦走入社会，可以说无时无刻不是为他人服务的：在外为公司做事，于内为家庭奉献。如果能以个人的知识、能力、成果给他人以有益的影响或带来直接的收益，则算是真正发挥了自身的价值。自然，要时时问自己：“为人谋而不忠乎？”面对承担的工作、参与或负责的项目，绝不是草草收场了事或者在发表一堆高谈阔论后半途而废，一定要拿出实实在在的成果并能够善始善终。为此，不论当前工作是否与个人的兴趣或长远发展相关，仅就这一份承诺与责任也要尽心竭力将事情做好。否则，留给他人与世界的就只是清高与无为，又如何在社会上立足？这也就是所谓的“虽执鞭之士，吾亦为之”。

然而，上面所说只是人的工具性一面，是与外界的接口。与之相对的，则是更应大力培养与展示的人性的一面。“君子不器”即是这个道理。在将本职工作做好的同时——至少做到没有明显的阙漏，要在有余力和闲暇的时候加紧提升自身实力以塑成专长，培养兴趣爱好以丰富个人生活。当然，一个人工作以后，相对于学校环境往往面临更多干扰，因而塑造与提升自我也变得愈加困难。但无论身处何时何处，人都不可能随心所欲改变他人和环境，相反地，一定要习惯在充满干扰的环境中，寻找与利用各种可能的机会，集中精力学习、训练，以达成最好的自己。

\begin{classic}
《论语·述而篇第七 7.12》 子曰：「富而可求也，虽执鞭之士，吾亦为之。如不可求，从吾所好。」
\end{classic}

\entrydate{2016.11.30}

\chapter{中年人的人生正途}

想必对于一些中年人来说，无时无刻不在咬噬心灵的苦闷与痛楚便是：他/她虽然通过长期的学习与积累，已经能够清晰地界定出和领悟到何谓至真、至善、至美；也掌握了基本的知识与技能，可以欣赏与把玩艺术的极致与科学的纯粹；然而却只能眼睁睁地目睹因为个人智力、能力、精力的不足，或者更多的是出于刻苦用功不够、缺乏坚持与训练，而尚不能亲身涉足于这样的境界，亲手做出具有极至品质的完美之作。

此时，狂者的进取将会导致在虚妄成就的自我麻痹中误入歧途。而狷者的有所不为则是怯懦与无能者逃避现实、向命运缴枪妥协的无奈之举。对此，恐怕惟有儒家所倡导的“知其不可而为之”的不屈不挠的精神，才是中年人今后人生的正途吧！

\begin{classic}
《论语·子路篇第十三 13.21》 子曰：「不得中行而与之，必也狂狷乎？狂者进取，狷者有所不为也。」

《论语·宪问篇第十四 14.38》 子路宿于石门。晨门曰：「奚自？」子路曰：「自孔氏。」曰：「是知其不可而为之者与？」
\end{classic}

\entrydate{2018.10.19}

\chapter{自信的来源}

自信的来源不是已有的学历、成果、成功，而是始终如一的刻苦、细心与虚心。

刻苦：只有刻苦学习形成深厚的积累，对于理知的理解与掌握才不至于似是而非，才能做到百分百的可靠。

\begin{classic}
《论语·里仁篇第四 4.6》 子曰：「我未见好仁者，恶不仁者。好仁者，无以尚之；恶不仁者，其为仁矣，不使不仁者加乎其身。有能一日用其力于仁矣乎？我未见力不足者。盖有之矣，我未之见也。」
\end{classic}

细心：失之毫厘，谬以千里。每一小步的推导、计算、编码都要“谨小慎微”，才能铸就严丝合缝般的精确无误。

\begin{classic}
《论语·学而篇第一 1.15》 子贡曰：「贫而无谄，富而无骄，何如？」子曰：「可也，未若贫而乐，富而好礼者也。」子贡曰：「诗云：『如切如磋，如琢如磨』，其斯之谓与？」子曰：「赐也，始可与言诗已矣，告诸往而知来者。」
\end{classic}

虚心：只要距离真正的底层逻辑哪怕尚存一线距离，就不能自以为自知，更不能不懂装懂。在科学的面前保持一颗谦卑的心态，才可能获得进步。

\begin{classic}
《论语·为政篇第二 2.17》 子曰：「由！诲女知之乎！知之为知之，不知为不知，是知也。」
\end{classic}

\entrydate{2019.09.22}

\chapter{夫我则不暇}

\begin{classic}
《论语·宪问篇第十四 14.29》 贡方人。子曰：「赐也，贤乎哉？夫我则不暇！」
\end{classic}

这句话在当下显得尤为重要：面对学术圈与职场中的各种负面消息与丑闻，我们普通人其实是没有时间去关注、评论或产生不必要的情绪的。长期以来，这个世界在人性与道德方面的崩坏已经很久了。只不过由于技术的不断发展大幅地改善了我们的物质生活，以至于我们生出了过多的从物质世界拓展到道德与精神世界的幻想。同时，表面上不断推陈出新的制度与游戏规则(但实质换汤不换药甚至每况愈下)，如考核指标、升职标准、头衔荣誉等等，也给了我们过多不切实际的希望。“我欲仁，斯仁至矣”，我们应该立刻以最正确的方式去做正确的事情，我们已经没有时间与垃圾人和垃圾的体制浪费时间了。

\entrydate{2025.05.01}

\chapter{不忘初心与君子不器}

\begin{classic}
《论语·为政篇第二 2.12》子曰：「君子不器。」
\end{classic}

虽然人们常提“不忘初心”，但这个“初心”究竟是什么呢？恐怕将其定义为终极目标与个人原则较为合适。可如果没有特殊的引导与机缘，受个人与家庭知识、能力、手艺、经历、体验的限制，一个人是无法从幼年开始就明确形成在持久行动的意义上而非仅存于理性头脑或言语中的终极目标与个人原则，亦即所谓的“初心”的。以这样的标准来看，初心虽难以形成，但却不妨碍构成其因素的孕育、萌发、培养、酝酿。人至中年，若能在较大的时间尺度上保存并重温曾经的记忆，则有助于寻找这些一直于潜移默化中影响我们人生道路、做事风格、关键抉择的基因或种子。

在我小时候上学前，由于父母是国企双职工，所以我经常是一个人在家，没法出去玩。当时的一件乐事便是将围棋棋子从盒子中一个个拿出来，严格保持间距、黑白交替地在沙发上排满。看到最终形成的规模“宏大”、整齐壮观的阵列，我心中会油然生出莫名的舒适感。想来，这便是我如今依然能够从数学与代码中获得美感的享受、一直乐在其中的重要原因。

我上小学后，在一次速算考试上，有一个同学在我之前率先交卷，但由于他的计算结果存在个别错误，所以最终我仍旧得到了第一名。这件事给我的印象很深、影响极大，以至于在我整个的研究与工作中，始终都将结果100\%的正确作为第一要务而非速度，“欲速则不达”。而这正是从事数值计算研究与开发工作的基本要求与学术修养。

大概从小学四年级开始，我在父母的指导与建议下开始有了自己的错题本。每次作业与考试中犯下的错误都会重新誉写在这个本子上，尤其是会总结出错的原因，并提示下次应该注意的事项。这样一来，我犯的错误逐渐变少，直至最后只会偶尔地出现。现在想来，这种习惯与素质对于我如今开发数值算法与软件是必不可缺的。首先，在我读过的几乎所有重要学术文献中都存在因为有意或无意的推导错误。这些错误虽然细微，但对于以此为依据做出自己的算法与软件却是致命的。因此，必须以百倍的谨慎、细心和信心发现并纠正这些错误。其次，如何将数学理论的正确性完全地转化为最终可运行代码的正确性也需要秉持同样的态度、品质与能力。毫无疑问，儿时在父母的教育下所形成的学习习惯为我打下了良好、坚实的基础，是日后能够以高标准把控工作流程与最终结果的保证。

除了上述错题本，在小学的每次考试中，由于我往往能够很快地完成试卷——大概只占整个考试时间的三分之一——所以会有大量的时间用于反复核对检查。而这需要极大的耐心与专注：一是要克制提前交卷的冲动；二是能够不厌其烦，重复几十遍做相同的推导与验证工作；三是为了避免掉入思维定势的陷阱，还得尝试以不同的角度与思路切入，来验证同一个问题。以现在的眼光来看，所有这些耐心、细心、专注恰恰是在算法调试与代码排错的枯燥过程中所极为需要的。

此外，在父母教育与个人性情的契合下，虽然我在小学与中学的每次考试中都能获得优异成绩，但我从未有过丝毫骄傲、自满的情绪与任何“邀功请赏”的行为。一方面，这保证了我在今后整个学习与研究历程中戒骄戒躁，避免因为一时的得意忘形而在学术成果、软件产品上出现无论是重大还是微小的疏漏。另一方面，更为重要的是，帮助我形成、确定并坚定了将学习本身作为目的、作为乐趣而绝非手段的基本思想与原则。想来，这也符合孔子所谓“君子不器”的精神。

以上将我的一些回忆与思考整理出来，才发现它们竟能连点成线，形成我如今学习与研究的原则、态度、品味、风格。想来“不忘初心”也便是这个意思吧！

\entrydate{2022.05.10}

\chapter{知天命}

\begin{classic}
《论语·为政篇第二 2.4》子曰：「吾十有五而志于学，三十而立，四十而不惑，五十而知天命，六十而耳顺，七十而从心所欲，不逾矩。」
\end{classic}

Even though I’m not as old as 50s, I’ve realized that I was born to do things which require great carefulness and rigorousness. This makes me reminiscent of the feeling of taking so many semester and Olympic exams during elementary school and middle school. In these exams, repeated checking every step of the solution, ensuring the correctness of every number and keeping the neatness of the calligraphy were all mandatory. Luckily, I had already got used to such kind of monotonous labor.

For years, I’ve never been satisfied with coarse and inexact derivations, which lack perfection and beauty. That’s why I initially started my college study in electrical engineering, but after years of perplexity and hard graft in darkness, I finally deviate from this area and devote my self in the field of mathematical analysis and numerical simulation.

Now, old habits, feelings and experiences from my youth just hit the spot of developing my mathematical software written in static language C++. Yes, both the mathematical theorem and programming language are of rigorous type, which just form my niche vocation.

\entrydate{2025.01.22}

\chapter{只问耕耘不问收获}

一直以来都喜欢大刘的《球状闪电》远远胜过《三体》。不仅是因为其中跌宕起伏的情节和细腻的人物描写，还因为在一开篇那个雨夜的场景中主人公爸爸说的一段话：

\begin{classic}
“其实，儿子，过一个美妙的人生并不难，听爸爸教你：你选一个公认的世界难题，最好是只用一张纸和一只铅笔的数学难题，比如歌德巴赫猜想或费尔马大定理什么的，或连纸笔都不要的纯自然哲学难题，比如宇宙的本源之类，投入全部身心钻研，只问耕耘不问收获，不知不觉的专注中，一辈子也就过去了。人们常说的寄托，也就是这么回事。或是相反，把挣钱作为惟一的目标，所有的时间都想着怎么挣，也不问挣来干什么用，到死的时候像葛朗台一样抱者一堆金币说：啊，真暖和啊……所以，美妙人生的关键在于你能迷上什么东西。”
\end{classic}

这里说的“只问耕耘不问收获”的最大难处还并非理论和技术上的障碍，而是在久无收获尤其是现实利益之后可能产生的自信心的动摇，对初衷、路线和终极目标的怀疑。但是，一个人的目光总是短浅的，即便是自以为理性的判断也都是构建在已有的知识和经验之上，因而是受限的。如果我们认定了积累与耕耘的过程本身就是意义，那么还有什么值得怀疑呢？更何况，若不照此埋头笃行，待到未来的某一天看到有人恰好循此路线将事情做成而自己却依旧在幻想中两手空空，且不是因为个人能力问题而是志向不够远大、信心不够坚定而导致了如此的天壤之别，那么我们对此难道不会有深深的遗憾与愧疚吗？

所以，一旦我们迷上了什么事情，便只管按照既定的方向踏踏实实地去做就是了，而不应将拿到了博士学位、评上了高级职称、找到了高收入的好工作等等量化指标的丰富或极大化当成是自我可以停止学习和成“器”的标志。而且“君子无终食之间违仁”，这种“痴迷”与时间、地点无关，与个人的身体与精神状态无关。其间任何无聊与痛苦的感受，甚至是放弃的想法，都是干扰的噪声，都是要予以立即忽略和摒弃的。

总之，我们每个人不应限制和辜负了自己无限的潜力和不断提升的可能，就是要“只问耕耘不问收获”地朝着真善美的方向走在虽然无比艰辛但却通向光明的路上。

\begin{classic}
《论语·雍也篇第六 6.12》 冉求曰：「非不说子之道，力不足也。」子曰：「力不足者，中道而废。今女画。」

《论语·为政篇第二 2.12》 子曰：「君子不器。」

《论语·里仁篇第四 4.5》 子曰：「富与贵，是人之所欲也，不以其道得之，不处也。贫与贱，是人之恶也，不以其道得之，不去也。君子去仁，恶乎成名。君子无终食之间违仁，造次必于是，颠沛必于是。」
\end{classic}

\entrydate{2018.05.13}
